
3.10
Soit @@M0@@ une @@M1@@-@@M2@@-algèbre augmentée. On définit sur @@M3@@ une filtration d'anneau de la manière suivante : pour tout @@M4@@, on désigne par @@M5@@ l'idéal engendré par les produits :
@@M6@@
@@M7@@ étant l'augmentation).

Supposons que @@M8@@ soit engendré en tant que groupe abélien par une famille d'éléments @@M9@@~; alors @@M10@@ est engendré comme @@M11@@-module par les produits :
@@M12@@
En effet, si @@M13@@ est un élément de @@M14@@ d'augmentation nulle et si on écrit @@M15@@, on a @@M16@@, donc @@M17@@~; par ailleurs, si @@M18@@, on peut écrire @@M19@@~; en utilisant les relations (3.4.1), on voit donc que les éléments de @@M20@@ sont combinaisons linéaires à coefficients dans @@M21@@ des éléments (3.10.2)~; comme ceux-ci sont dans @@M22@@, on obtient le résultat voulu.

La filtration ainsi définie est une filtration d'anneau~; elle est décroissante, et on a :
@@M23@@
De plus, elle est fonctorielle par rapport à @@M24@@ : tout @@M25@@-homomorphisme de @@M26@@-@@M27@@-algèbres augmentées est compatible aux filtrations. Quand nous considérerons un @@M28@@-anneau augmenté comme muni d'une filtration, il s'agira, sauf mention expresse du contraire, de la filtration ci-dessus.

